\cleardoublepage
\phantomsection
\begin{planbox}[แผนบริหารการสอนประจำบทที่ 9\quad ทัศนศาสตร์พื้นฐาน การสะท้อน และการหักเหของแสง (สัปดาห์ที่ 10)]
\noindent\textbf{1. หัวข้อเนื้อหาประจำบท (บรรยาย 2 ชั่วโมง และปฏิบัติการ 3 ชั่วโมง)}
\begin{itemize}
    \item 9.1 ธรรมชาติของแสง สเปกตรัมคลื่นแม่เหล็กไฟฟ้า และดัชนีหักเหของตัวกลาง ($n = c/v$)
    \item 9.2 กฎการสะท้อนของแสงและกฎการหักเหของสเนลล์ในบริบทการส่องสไลด์จุลชีววิทยา
    \item 9.3 การประยุกต์ใช้เครื่องวัดดัชนีหักเห (Refractometer) วัดความเข้มข้นน้ำตาลในการหมักด้วยจุลินทรีย์
    \item 9.4 ปฏิบัติการที่ 9 การศึกษาทางเดินรังสีแสงสะท้อน แสงหักเห และการวัดค่าบริกซ์ (Brix) ของอาหารหมัก
\end{itemize}

\noindent\textbf{2. วัตถุประสงค์เชิงพฤติกรรมและผลลัพธ์การเรียนรู้รายวิชา (CLOs ที่เชื่อมโยง PLOs)}
\begin{enumerate}
    \item อธิบายความหมายของดัชนีหักเห ($n$) และทิศทางการเบนของลำแสงเมื่อเดินทางผ่านตัวกลางต่างชนิดกันได้ \textbf{(CLO1 เชื่อมโยง PLO2)}
    \item อธิบายหลักการทำงานของเครื่อง Hand Refractometer ในการวัดปริมาณน้ำตาล ($^\circ\text{Brix}$) ของกระบวนการหมักไวน์หรืออาหารเลี้ยงเชื้อได้ \textbf{(CLO1, CLO2 เชื่อมโยง PLO2, PLO3)}
    \item ปฏิบัติการทดลองลากแนวรังสีแสงตกกระทบและรังสีหักเหอย่างละเอียดรอบคอบ \textbf{(CLO3, CLO4 เชื่อมโยง PLO1, PLO5, PLO6)}
\end{enumerate}

\noindent\textbf{3. กิจกรรมการจัดการเรียนรู้เชิงรุก (Active Learning \& OBE)}
\begin{itemize}
    \item \textbf{Hands-on Optical Ray Tracing}\quad ใช้กล่องกำเนิดลำแสงเลเซอร์ฉายผ่านแท่งแก้วปริซึมและสารละลายน้ำตาลความเข้มข้นต่างกัน พร้อมจำกฎง่ายๆ โดยไม่ต้องกดเครื่องคิดเลขตรีโกณมิติว่า ``\textbf{น้อยไปมาก เบนเข้าหาเส้นแนวฉาก / มากไปน้อย เบนออกจากเส้นแนวฉาก}''
\end{itemize}

\noindent\textbf{4. สื่อการเรียนรู้และแหล่งเรียนรู้}
\begin{itemize}
    \item เอกสารคำสอนบทที่ 9 ชุดกล่องแสงเลเซอร์ แท่งพลาสติกใส และเครื่องวัดความหวาน Hand Refractometer
\end{itemize}

\noindent\textbf{5. การวัดและประเมินผลประจำบท}
\begin{itemize}
    \item คำถามท้ายคาบเรื่องดัชนีหักเหและทิศทางการเบนของแสง \textbf{(CLO1, CLO2)}
    \item รายงานปฏิบัติการที่ 9 \textbf{(CLO3, CLO4)}
\end{itemize}
\end{planbox}
\clearpage

\chapter{ทัศนศาสตร์พื้นฐาน การสะท้อน และการหักเหของแสง}
\index{ทัศนศาสตร์}\index{ดัชนีหักเห}\index{กฎของสเนลล์}\index{เครื่องวัดความหวาน Refractometer}

\section{ธรรมชาติของแสงและดัชนีหักเหของตัวกลางทางชีวภาพ}

\deepskinfig{figures/sec_3d/ch09_s1.jpg}{เปรียบเทียบค่าดัชนีหักเห ($n = c/v$) ของอากาศ น้ำ ไซโทพลาซึมของเซลล์สด กระจกสไลด์ และน้ำมันช่วยส่อง (Immersion Oil, $n = 1.515$)}

แสงที่ตามนุษย์มองเห็น (\textbf{Visible Light}) เป็นคลื่นแม่เหล็กไฟฟ้าที่มีช่วงความยาวคลื่น ($\lambda$) อยู่ระหว่าง $400\text{ nm}$ (แสงสีม่วง) ถึง $700\text{ nm}$ (แสงสีแดง) ในสุญญากาศหรือในอากาศ แสงเดินทางด้วยอัตราเร็วสูงสุดคงที่คือ $c \approx 3.0 \times 10^8\text{ m/s}$ แต่เมื่อแสงเดินทางเข้าสู่ตัวกลางโปร่งใสอื่นๆ เช่น น้ำ กลีเซอรีน กระจกสไลด์ หรือน้ำมันช่วยส่อง (\textbf{Immersion Oil}) อัตราเร็วของแสง ($v$) จะลดลง อัตราส่วนระหว่างอัตราเร็วแสงในสุญญากาศต่ออัตราเร็วแสงในตัวกลางนั้น เรียกว่า \textbf{ดัชนีหักเห} (\textbf{Refractive Index}, $n$) ดังสมการที่ (9.1)

\begin{equation}
    n = \frac{c}{v}
\end{equation}

ตัวกลางที่มีความหนาแน่นเชิงแสงมากจะมีค่าดัชนีหักเห $n$ สูง และแสงจะเดินทางผ่านได้ช้าลง ค่าดัชนีหักเหของตัวกลางที่เกี่ยวข้องกับการส่องกล้องจุลทรรศน์แสดงไว้ใน\mbox{ภาพที่} 9.1 และ\mbox{ตารางที่} 9.1

\begin{table}[H]
\centering
\caption{ค่าดัชนีหักเห ($n$) ของตัวกลางสำคัญในงานกล้องจุลทรรศน์และจุลชีววิทยา}
\small
\begin{tabularx}{\textwidth}{>{\bfseries}p{4.5cm} c >{\raggedright\arraybackslash}X}
\toprule
ชนิดของตัวกลาง & ดัชนีหักเห ($n$) & ความสำคัญในงานจุลชีววิทยา \\
\midrule
อากาศ (Air) & $1.000$ & ช่องว่างระหว่างแผ่นสไลด์กับเลนส์ใกล้วัตถุกำลังขยายต่ำ (4x, 10x, 40x) \\
น้ำบริสุทธิ์ (Water) & $1.333$ & ตัวกลางหลักภายในไซโทพลาซึมของเซลล์สด (ทำให้เซลล์สดดูใสเมื่อไม่ย้อมสี) \\
กระจกสไลด์และแผ่นปิดสไลด์ & $1.515$ & วัสดุมาตรฐานของแผ่นสไลด์แก้วและเลนส์ใกล้วัตถุของกล้องจุลทรรศน์ \\
น้ำมันช่วยส่อง (Immersion Oil) & $1.515$ & มีค่า $n$ เท่ากับกระจกสไลด์ทุกประการ ใช้หยดเชื่อมเลนส์หัวน้ำมัน 100x เพื่อป้องกันแสงหักเหหนี \\
\bottomrule
\end{tabularx}
\end{table}

\section{กฎการสะท้อนและกฎการหักเหของแสงฉบับไม่ต้องจำตรีโกณมิติ}

\deepskinfig{figures/sec_3d/ch09_s2.jpg}{การเบนของลำแสงเลเซอร์ตามกฎการหักเหของสเนลล์เมื่อเดินทางผ่านรอยต่อระหว่างตัวกลางที่มีค่าดัชนีหักเหต่างกัน}

เมื่อลำแสงตกกระทบผิวรอยต่อระหว่างตัวกลางสองชนิดโดยทำมุมตกกระทบ $\theta_1$ กับ\textbf{เส้นแนวฉาก} (\textbf{Normal Line} คือเส้นประที่ลากตั้งฉากกับผิวรอยต่อ) แสงส่วนหนึ่งจะสะท้อนกลับโดยมี\textbf{มุมตกกระทบเท่ากับมุมสะท้อนเสมอ} ($\theta_{\text{incident}} = \theta_{\text{reflected}}$) และแสงอีกส่วนหนึ่งจะทะลุผ่านเข้าสู่ตัวกลางที่สองโดยเกิดการเบนทิศทางตาม\textbf{กฎของสเนลล์} (\textbf{Snell's Law}) ในสมการที่ (9.2)

\begin{equation}
    n_1 \sin\theta_1 = n_2 \sin\theta_2
\end{equation}

สำหรับนักศึกษาจุลชีววิทยา เราสามารถสรุปทิศทางการหักเหของแสงเพื่อนำไปใช้วิเคราะห์กล้องจุลทรรศน์ได้ทันทีโดยไม่ต้องกดเครื่องคิดเลขหาค่าไซน์ ($\sin$) ดังแสดงใน\mbox{ภาพที่} 9.2 และ\mbox{ภาพที่} 9.3

\begin{figure}[H]
\centering
\resizebox{0.86\textwidth}{!}{%
\begin{tikzpicture}[font=\sffamily, >=Stealth]
    % Left: Air (n=1.0) to Glass (n=1.515) -> Bend TOWARD Normal
    \begin{scope}[shift={(0,0)}]
        \fill[boxBlueBg] (0,0) rectangle (6.4,2.2);
        \fill[white] (0,2.2) rectangle (6.4,4.4);
        \draw[physRoyalBlue, line width=1.4pt] (0,2.2) rectangle (6.4,4.4);
        \draw[physRoyalBlue, line width=1.4pt] (0,0) rectangle (6.4,2.2);
        \draw[dashed, physSlate, line width=1.1pt] (3.2,0.2) -- (3.2,4.2);

        \draw[->, crimsonAlert, line width=1.8pt] (1.2,4.0) -- (3.2,2.2);
        \draw[->, crimsonAlert, line width=1.8pt] (3.2,2.2) -- (4.1,0.3);

        \node[font=\footnotesize\bfseries, color=physSlate] at (1.3,2.5) {อากาศ ($n_1 = 1.0$)};
        \node[font=\footnotesize\bfseries, color=physRoyalBlue] at (1.5,1.8) {กระจกสไลด์ ($n_2 = 1.515$)};
        \node[font=\small\bfseries, color=physRoyalBlue] at (3.2,-0.45) {(ก) $n$ น้อย $\rightarrow$ $n$ มาก (เบนเข้าหาเส้นฉาก)};
    \end{scope}

    % Right: Glass (n=1.515) to Air (n=1.0) -> Bend AWAY from Normal
    \begin{scope}[shift={(7.6,0)}]
        \fill[white] (0,2.2) rectangle (6.4,4.4);
        \fill[boxBlueBg] (0,0) rectangle (6.4,2.2);
        \draw[crimsonAlert, line width=1.4pt] (0,0) rectangle (6.4,4.4);
        \draw[crimsonAlert, line width=1.2pt] (0,2.2) -- (6.4,2.2);
        \draw[dashed, physSlate, line width=1.1pt] (3.2,0.2) -- (3.2,4.2);

        \draw[->, bioEmerald, line width=1.8pt] (2.2,0.3) -- (3.2,2.2);
        \draw[->, crimsonAlert, line width=1.8pt] (3.2,2.2) -- (5.8,3.6);

        \node[font=\footnotesize\bfseries, color=physSlate] at (1.4,2.5) {อากาศ ($n_2 = 1.0$)};
        \node[font=\footnotesize\bfseries, color=physRoyalBlue] at (1.6,1.8) {กระจกสไลด์ ($n_1 = 1.515$)};
        \node[font=\small\bfseries, color=crimsonAlert] at (3.2,-0.45) {(ข) $n$ มาก $\rightarrow$ $n$ น้อย (เบนหนีออกจากเส้นฉาก)};
    \end{scope}
\end{tikzpicture}%
}
\caption{ทิศทางการเบนของลำแสงเมื่อเดินทางผ่านรอยต่อระหว่างอากาศกับกระจกสไลด์}
\end{figure}

\vspace{4pt}
\arucomarkerbadge{8}{9}{การหักเหของสเนลล์และเครื่องวัดค่าบริกซ์การหมักยีสต์ Refractometer}{แสดงการเบนของลำแสงเลเซอร์เมื่อผ่านน้ำหมักที่มีความเข้มข้นน้ำตาลต่างกัน และสเกลเส้นแบ่งเขตมืด--สว่างในหน่วยองศาบริกซ์ ($^\circ\text{Brix}$)}{มุมตกกระทบ $\theta_1$ ($^\circ$), ค่าความหวานน้ำหมัก ($0\text{--}30\,^\circ\text{Brix}$), เวลาการหมักยีสต์}


\begin{worked}[การคำนวณอัตราเร็วของแสงเมื่อเดินทางผ่านแผ่นสไลด์แก้วและน้ำ]
\textbf{โจทย์ปัญหา}\quad แสงเดินทางในสุญญากาศด้วยอัตราเร็ว $c = 3.0 \times 10^8\text{ m/s}$ เมื่อลำแสงจากหลอดไฟใต้กล้องจุลทรรศน์เดินทางผ่านหยดน้ำบนสไลด์ซึ่งมีดัชนีหักเห $n = 1.33$ อัตราเร็วของแสงในหยดน้ำ ($v$) จะมีค่าเท่าใด

\textbf{วิธีทำทีละขั้นตอนฉบับเข้าใจง่าย}
\begin{enumerate}
    \item \textbf{ขั้นที่ 1 ใช้สามเหลี่ยมสูตรดัชนีหักเห $n = c / v$ ปิดตัวแปร $v$ จะได้ $v = c / n$}
    \[
        v = \frac{c}{n}
    \]
    \item \textbf{ขั้นที่ 2 แทนค่าตัวเลขแล้วหารกัน}
    \[
        v = \frac{3.0 \times 10^8\text{ m/s}}{1.33} \approx 2.26 \times 10^8\text{ m/s}
    \]
\end{enumerate}
\textbf{คำตอบ}\quad อัตราเร็วของแสงขณะเดินทางผ่านหยดน้ำบนสไลด์มีค่าเท่ากับ $2.26 \times 10^8\text{ m/s}$
\end{worked}

\section{ปฏิบัติการประจำสัปดาห์ที่ 10}

\deepskinfig{figures/sec_3d/ch09_s3.jpg}{ชุดปฏิบัติการที่ 9 สำหรับวัดมุมตกกระทบ--มุมหักเหผ่านแท่งอะคริลิกใส และการวัดความหวานน้ำหมักยีสต์ ($^\circ\text{Brix}$) ด้วยเครื่อง Hand Refractometer}

\begin{labbox}[ปฏิบัติการที่ 9\quad การสะท้อนและการหักเหของแสงผ่านตัวกลางและสารละลาย (3 ชั่วโมง)]
\noindent\textbf{ผลลัพธ์การเรียนรู้ที่ประเมิน}\quad CLO1 (การสะท้อนและการหักเหของแสง), CLO3 (ทักษะการทดลองทางแสง), CLO4 (รายงานกลุ่ม)

\vspace{4pt}
\noindent\textbf{1. วัสดุและอุปกรณ์ (ดังแสดงใน\mbox{ภาพที่} 9.4)}
\begin{itemize}
    \item ชุดกล่องกำเนิดลำแสงเลเซอร์, กระจกเงาระนาบ, แท่งอะคริลิกใสครึ่งวงกลม, กระดาษกราฟวงกลมวัดมุม (Protractor)
    \item เครื่องวัดดัชนีหักเหและความหวาน (Hand Refractometer) และสารละลายน้ำตาลกลูโคสความเข้มข้น $5\%, 10\%$ และ $20\%$
\end{itemize}

\noindent\textbf{2. ขั้นตอนการทดลอง}
\begin{enumerate}
    \item \textbf{ตอนที่ 1 การตรวจสอบกฎการหักเหของแสง}\quad วางแท่งอะคริลิกใสบนกระดาษวัดมุม ฉายลำแสงเลเซอร์จากอากาศเข้าสู่อะคริลิกด้วยมุมตกกระทบ $\theta_1 = 30^\circ$ และ $45^\circ$ อ่านค่ามุมหักเห $\theta_2$ และสังเกตว่าลำแสงเบนเข้าหาเส้นแนวฉากจริงหรือไม่
    \item \textbf{ตอนที่ 2 การใช้เครื่อง Refractometer ติดตามความเข้มข้นน้ำตาลสำหรับการหมักจุลินทรีย์}\quad หยดน้ำกลั่นปรับตั้งศูนย์เครื่อง Hand Refractometer จากนั้นหยดสารละลายน้ำตาลกลูโคส $5\%, 10\%$ และ $20\%$ อ่านค่าองศาบริกซ์ ($^\circ\text{Brix}$) ซึ่งแปรผันตรงกับดัชนีหักเหของสารละลาย
\end{enumerate}

\noindent\textbf{3. ตารางบันทึกผลการทดลอง}
\begin{center}
\footnotesize\setlength{\tabcolsep}{3.5pt}
\begin{tabularx}{\linewidth}{>{\raggedright\arraybackslash}X c c >{\raggedright\arraybackslash}X}
\toprule
\textbf{ตัวอย่างสารละลายและการทดลอง} & \textbf{มุมตกกระทบ $\theta_1$ หรือความเข้มข้น} & \textbf{มุมหักเห $\theta_2$ หรือค่า $^\circ\text{Brix}$} & \textbf{สรุปทิศทางการเบนของแสง} \\
\midrule
แสงผ่านอากาศสู่อะคริลิก ครั้งที่ 1 & $\theta_1 = 30^\circ$ & $\theta_2 = $ ............ & \dotfill \\
แสงผ่านอากาศสู่อะคริลิก ครั้งที่ 2 & $\theta_1 = 45^\circ$ & $\theta_2 = $ ............ & \dotfill \\
สารละลายน้ำตาลกลูโคสหมักยีสต์ & $10\%$ และ $20\%$ & ............ / ............ $^\circ\text{Brix}$ & \dotfill \\
\bottomrule
\end{tabularx}
\end{center}
\end{labbox}

\section*{คำถามท้ายบทที่ 9 (ทบทวนแนวคิดและการคำนวณเบื้องต้น)}
\begin{enumerate}
    \item เพราะเหตุใดการส่องดูเซลล์แบคทีเรียที่มีชีวิตในหยดน้ำ (Wet Mount) โดยไม่ย้อมสี จึงมองเห็นขอบเขตเซลล์ได้ยากมาก จงอธิบายโดยเปรียบเทียบดัชนีหักเหของไซโทพลาซึมกับน้ำรอบเซลล์
    \item เมื่อแสงเดินทางออกจากแผ่นสไลด์แก้ว ($n = 1.515$) ออกสู่ชั้นอากาศ ($n = 1.000$) ลำแสงจะเบนเข้าหาหรือเบนหนีออกจากเส้นแนวฉาก
\end{enumerate}
